\documentclass[11pt,a4paper]{article}

% --- Codificacion e idioma -------------------------------------------------
% NOTA: se evita deliberadamente el paquete babel con el idioma espanol,
% ya que requiere el paquete adicional "texlive-lang-spanish" que no
% siempre esta instalado por defecto (riesgo de romper la compilacion en
% una maquina distinta). Los textos fijos en espanol (indice, tablas,
% referencias) se traducen manualmente mas abajo; el contenido del
% documento ya esta redactado en espanol de todas formas.
\usepackage[utf8]{inputenc}
\usepackage[T1]{fontenc}

% --- Layout y utilidades ----------------------------------------------------
\usepackage[margin=2.5cm]{geometry}
\usepackage{graphicx}
\usepackage{booktabs}
\usepackage{array}
\usepackage{xcolor}
\usepackage{listings}
\usepackage{hyperref}
\usepackage{enumitem}
\usepackage{fancyhdr}

\hypersetup{
    colorlinks = true,
    linkcolor  = black,
    urlcolor   = blue,
    citecolor  = black,
    pdftitle   = {Informe Tecnico - Sistema de Inventario Mercado Municipal de Quevedo},
    pdfauthor  = {Estudiante UTEQ}
}

\pagestyle{fancy}
\fancyhf{}
\fancyhead[L]{Aplicaciones Web -- Ingenieria de Software -- UTEQ}
\fancyhead[R]{Examen Parcial -- Informe Tecnico}
\fancyfoot[C]{\thepage}
\setlength{\headheight}{14pt}

% Traduccion manual de textos fijos (sin depender de babel/spanish)
\renewcommand{\contentsname}{Indice}
\renewcommand{\tablename}{Tabla}
\renewcommand{\refname}{Referencias}

% --- Estilo para bloques de codigo / SQL / JSON -----------------------------
\definecolor{codebg}{RGB}{244,244,244}
\definecolor{codeborder}{RGB}{200,200,200}

\lstdefinestyle{codigo}{
    backgroundcolor=\color{codebg},
    basicstyle=\ttfamily\footnotesize,
    breaklines=true,
    frame=single,
    rulecolor=\color{codeborder},
    numbers=none,
    tabsize=2,
    columns=fullflexible,
    keepspaces=true
}
\lstset{style=codigo}

\title{\Large \textbf{Sistema de Gestion de Inventario} \\ \large Mercado Municipal de Quevedo}
\author{}
\date{}

\begin{document}

% =============================================================================
% CARATULA
% =============================================================================
\begin{titlepage}
    \centering
    \vspace*{1cm}

    {\large Universidad Tecnica Estatal de Quevedo\par}
    {\large Facultad de Ciencias de la Computacion y Diseno Digital\par}
    {\large Carrera de Ingenieria de Software (Redisenio)\par}
    \vspace{0.5cm}
    {\normalsize Aplicaciones Web -- Quinto Nivel\par}
    {\normalsize Periodo Academico Presencial 2026-2027 (PPA)\par}

    \vspace{2cm}
    \rule{\linewidth}{0.5pt}
    \vspace{0.5cm}
    {\Huge \textbf{Informe Tecnico}\par}
    \vspace{0.3cm}
    {\Large Sistema de Gestion de Inventario\par}
    {\Large Mercado Municipal de Quevedo\par}
    \vspace{0.3cm}
    {\large Examen Parcial -- Parte Practica\par}
    \rule{\linewidth}{0.5pt}

    \vspace{2cm}
    \begin{tabular}{ll}
        \textbf{Apellidos y nombres:} & [MARISCAL CABRERA JAIME JOSUÉ] \\[4pt]
        \textbf{Cedula:}              & [1250710835] \\[4pt]
        \textbf{Paralelo:}            & [A] \\[4pt]
        \textbf{Asignatura:}          & Aplicaciones Web \\[4pt]
        \textbf{Fecha del examen:}    & [24/7/2026] \\[4pt]
        \textbf{Docente:}             & Dr. Gleiston Ciceron Guerrero Ulloa, Ph.D. \\
    \end{tabular}

    \vspace{2cm}
    \fbox{
        \begin{minipage}{0.9\linewidth}
            \centering
            \vspace{0.3cm}
            \textbf{Repositorio del proyecto (URL en una sola linea):}\\[6pt]
            \url{}
            \vspace{0.3cm}
        \end{minipage}
    }

    \vfill
    {\small Este documento sirve como caratula de identificacion para el
    criterio de piso P1 y como portada del informe tecnico (criterio P2).
    \textbf{Reemplazar el enlace de arriba por la URL real del repositorio
    publico antes de compilar la version final.}\par}

\end{titlepage}

\tableofcontents
\newpage

% =============================================================================
\section{Introduccion}
% =============================================================================

El Municipio de Quevedo requiere digitalizar el control de productos del
mercado central. Este informe documenta el diseno, la implementacion y la
verificacion de un servicio web REST construido con Spring Boot~3 sobre
Java~21~LTS, que expone el recurso \texttt{Producto} bajo un contrato JSON
uniforme, valida la entrada en el servidor, aplica borrado logico, integra
el patron cache-aside con Redis y protege sus operaciones con autenticacion
y autorizacion basadas en JSON Web Token (JWT)~\cite{rfc7519}.

El trabajo cubre los cinco requisitos funcionales exigidos por el
enunciado del examen parcial, y se enmarca en los conocimientos de
construccion de software y calidad descritos por el SWEBOK
Guide~\cite{swebok2024}, como prerrequisito directo del Proyecto Fin de
Curso:

\begin{enumerate}[label=R\arabic*.]
    \item Listado paginado de productos activos, con envoltura de respuesta uniforme.
    \item Creacion de productos con validacion de entrada en el servidor.
    \item Eliminacion logica (soft delete) de productos.
    \item Patron cache-aside con Redis sobre el listado de productos.
    \item Seguridad basada en JWT con autorizacion por roles.
\end{enumerate}

% =============================================================================
\section{Arquitectura y decisiones de diseno}
% =============================================================================

\subsection{Arquitectura por capas}

El proyecto sigue una arquitectura por capas clasica, con responsabilidades
claramente separadas y una unica direccion de dependencia (controlador
$\rightarrow$ servicio $\rightarrow$ repositorio), tal como recomiendan
Martin~\cite{martin2017} y Richards y Ford~\cite{richards2020}, y
siguiendo las convenciones idiomaticas de Spring Boot descritas por
Walls~\cite{walls2022} (inyeccion de dependencias por constructor,
starters de configuracion automatica, DTOs separados de las entidades
JPA):

\begin{itemize}
    \item \textbf{Controlador} (\texttt{ProductoController}, \texttt{AuthController}):
          expone los endpoints REST, delega toda la logica de negocio al
          servicio y no contiene acceso a datos.
    \item \textbf{Servicio} (\texttt{ProductoServiceImpl}): contiene las
          reglas de negocio (paginacion, soft delete, cache-aside) y
          orquesta al repositorio.
    \item \textbf{Repositorio} (\texttt{ProductoRepository},
          \texttt{UsuarioRepository}): unico punto de acceso a la base de
          datos, mediante Spring Data JPA. Ninguna consulta se concatena
          manualmente con datos del usuario (se usan siempre parametros
          nombrados o derivados de metodo), evitando el riesgo de inyeccion
          SQL descrito en OWASP Top 10:2021 -- A03:2021 Injection~\cite{owasp2021}.
\end{itemize}

\subsection{Estilo arquitectonico REST}

La API sigue los principios de estilo arquitectonico REST descritos por
Fielding~\cite{fielding2000}: recursos identificados por URI
(\texttt{/api/v1/productos}), operaciones mapeadas a verbos HTTP (GET,
POST, DELETE) y ausencia de estado de sesion en el servidor (autenticacion
stateless mediante JWT).

% =============================================================================
\section{Modelo de datos}
% =============================================================================

El esquema se define explicitamente en \texttt{db/schema.sql} (tambien
incluido en \texttt{src/main/resources/schema.sql} para que Spring Boot lo
ejecute automaticamente al arrancar). Hibernate se configura en modo
\texttt{validate}: no genera el esquema, solo valida que las entidades
JPA coincidan con las tablas reales.

\subsection{Tabla \texttt{productos}}

\begin{lstlisting}[language=SQL]
CREATE TABLE IF NOT EXISTS productos (
    id         BIGSERIAL PRIMARY KEY,
    nombre     VARCHAR(100)   NOT NULL,
    categoria  VARCHAR(50)    NOT NULL,
    stock      INTEGER        NOT NULL CHECK (stock >= 0),
    precio     DECIMAL(10, 2) NOT NULL CHECK (precio >= 0.01),
    activo     BOOLEAN        NOT NULL DEFAULT TRUE,
    creado_en  TIMESTAMPTZ    NOT NULL DEFAULT now()
);
\end{lstlisting}

Las restricciones \texttt{CHECK} se aplican a nivel de base de datos como
defensa adicional, ademas de la validacion de Bean Validation en el DTO de
entrada (seccion~\ref{sec:validacion}).

\subsection{Tablas \texttt{usuarios} y \texttt{usuario\_roles}}

Los usuarios de autenticacion tambien se persisten en PostgreSQL (no estan
hardcodeados en el codigo Java), con la contrasena cifrada mediante BCrypt:

\begin{lstlisting}[language=SQL]
CREATE TABLE IF NOT EXISTS usuarios (
    id       BIGSERIAL PRIMARY KEY,
    username VARCHAR(50)  NOT NULL UNIQUE,
    password VARCHAR(100) NOT NULL
);

CREATE TABLE IF NOT EXISTS usuario_roles (
    usuario_id BIGINT      NOT NULL REFERENCES usuarios(id) ON DELETE CASCADE,
    rol        VARCHAR(20) NOT NULL,
    PRIMARY KEY (usuario_id, rol)
);
\end{lstlisting}

\texttt{db/seed.sql} inserta dos usuarios de prueba con hash BCrypt real
(generado con costo 10): \texttt{admin} (roles \texttt{ROLE\_ADMIN} y
\texttt{ROLE\_USER}) y \texttt{usuario} (solo \texttt{ROLE\_USER}).

% =============================================================================
\section{Contrato de la API}
% =============================================================================

Todas las respuestas comparten la misma envoltura JSON, inspirada en el
espiritu de uniformidad de \emph{Problem Details for HTTP
APIs}~\cite{rfc7807}, aunque adaptada al formato especifico exigido por el
enunciado:

\begin{lstlisting}
{ "success": bool, "data": T, "message": string,
  "meta": object | ausente, "errors": array | ausente }
\end{lstlisting}

\texttt{data} y \texttt{message} siempre se serializan (incluso cuando
\texttt{data} es \texttt{null}); \texttt{meta} solo aparece en el listado
paginado y \texttt{errors} solo aparece en las respuestas 400 de
validacion.

\subsection{Listado paginado -- 200 OK}
\begin{lstlisting}
GET /api/v1/productos?page=0&size=10&sort=nombre,asc

{
  "success": true,
  "data": [ { "id": 1, "nombre": "Arroz flor", "categoria": "Granos",
              "stock": 120, "precio": 0.95, "activo": true,
              "creadoEn": "2026-07-24T09:15:00-05:00" } ],
  "message": "Productos obtenidos exitosamente",
  "meta": { "page": 0, "size": 10, "totalElements": 57, "totalPages": 6 }
}
\end{lstlisting}

\subsection{Creacion invalida -- 400 Bad Request} \label{sec:validacion}
\begin{lstlisting}
{
  "success": false,
  "data": null,
  "message": "Error de validacion",
  "errors": [
    { "field": "nombre", "message": "no debe estar vacio" },
    { "field": "precio", "message": "debe ser mayor o igual a 0.01" }
  ]
}
\end{lstlisting}

La validacion se implementa con \texttt{@Valid} sobre el DTO de entrada
(\texttt{ProductoRequestDTO}), anotado con \texttt{@NotBlank},
\texttt{@Min(0)} y \texttt{@DecimalMin("0.01")}. Un
\texttt{@RestControllerAdvice} (\texttt{GlobalExceptionHandler}) captura
\texttt{MethodArgumentNotValidException} y construye la lista de errores
por campo.

% =============================================================================
\section{Seguridad: autenticacion y autorizacion con JWT}
% =============================================================================

La API se protege con JSON Web Token~\cite{rfc7519}. El flujo es el
siguiente:

\begin{enumerate}
    \item \texttt{POST /api/v1/auth/login} autentica usuario/contrasena
          contra la base de datos (\texttt{UserDetailsServiceImpl} +
          \texttt{DaoAuthenticationProvider} + \texttt{BCryptPasswordEncoder})
          y devuelve un token firmado con HS256, con los roles del usuario
          como \emph{claim}.
    \item \texttt{JwtAuthenticationFilter} valida el token en cada peticion
          y reconstruye la autenticacion en el \texttt{SecurityContext}.
    \item \texttt{SecurityConfig} exige \texttt{ROLE\_USER} para el GET y
          \texttt{ROLE\_ADMIN} para el POST y el DELETE.
\end{enumerate}

Los casos de error se resuelven con manejadores dedicados, evitando
confundir autenticacion con autorizacion (un error frecuente senalado en
OWASP Top 10:2021~\cite{owasp2021}):

\begin{itemize}
    \item \textbf{401 Unauthorized} (\texttt{JwtAuthenticationEntryPoint}):
          la peticion no trae token, o el token es invalido/expirado.
    \item \textbf{403 Forbidden} (\texttt{CustomAccessDeniedHandler}): el
          token es valido, pero el rol del usuario no alcanza para la
          operacion solicitada.
\end{itemize}

El secreto de firma (\texttt{JWT\_SECRET}) \textbf{no se encuentra
embebido en el repositorio}: se lee exclusivamente de una variable de
entorno, provista mediante un archivo \texttt{.env} (no versionado) a
partir de la plantilla \texttt{.env.example} con un valor ficticio.

% =============================================================================
\section{Cache-aside con Redis}
% =============================================================================

El listado de productos aplica el patron cache-aside descrito por
Kleppmann~\cite{kleppmann2017}:

\begin{itemize}
    \item \texttt{@Cacheable(value = "productos", key = ...)} sobre
          \texttt{ProductoServiceImpl.listar(...)}: en un \emph{cache miss}
          consulta PostgreSQL, arma la respuesta paginada y la almacena en
          Redis; en un \emph{cache hit} la sirve directamente desde Redis
          sin tocar la base de datos.
    \item \texttt{@CacheEvict(value = "productos", allEntries = true)}
          sobre \texttt{crear(...)} y \texttt{eliminar(...)}: cualquier
          escritura invalida todo el cache de listados, para no servir
          datos obsoletos tras un alta o una baja.
    \item El TTL del espacio de nombres \texttt{productos} se configura en
          10 minutos (\texttt{application.yml}), como limite superior de
          obsolescencia si ninguna escritura ocurre.
\end{itemize}

La clave de cache (\texttt{page-size-sort}) evita colisiones entre
distintas combinaciones de paginacion y orden, ya que cada una es un
resultado distinto que debe cachearse por separado.

% =============================================================================
\section{Reproducibilidad}
% =============================================================================

El repositorio se puede clonar y ejecutar en una maquina limpia con un
unico comando, sin configuracion manual adicional (mas alla de copiar la
plantilla de variables de entorno):

\begin{lstlisting}[language=bash]
git clone https://github.com/USUARIO/inventario-mercado-APELLIDO.git
cd inventario-mercado-APELLIDO
cp .env.example .env
docker compose up -d --build
\end{lstlisting}

Este comando levanta tres contenedores: \texttt{postgres} (con
\emph{healthcheck}), \texttt{redis} (con \emph{healthcheck}) y \texttt{app}
(construido desde el \texttt{Dockerfile} multi-stage del proyecto, que
compila con Maven en la primera etapa y ejecuta con un JRE 21 liviano en
la segunda). El servicio \texttt{app} espera a que \texttt{postgres} y
\texttt{redis} esten saludables (\texttt{condition: service\_healthy})
antes de arrancar.

Al iniciar, la aplicacion ejecuta automaticamente \texttt{schema.sql} y
\texttt{data.sql} (identicos a \texttt{db/schema.sql} y
\texttt{db/seed.sql}), dejando la base de datos lista con la tabla
\texttt{productos} vacia y los dos usuarios de prueba ya sembrados.

% =============================================================================
\section{Pruebas realizadas}
% =============================================================================

La coleccion \texttt{docs/requests.http} ejercita los cinco requisitos y
los cuatro codigos de estado de error exigidos. Siguiendo el principio de
reproducibilidad de la experimentacion en ingenieria de software descrito
por Wohlin et al.~\cite{wohlin2012}, cada caso de prueba se documenta con
su entrada, su codigo HTTP esperado y su criterio de verificacion, de modo
que un tercero pueda repetirlo exactamente. La tabla~\ref{tab:pruebas}
resume los casos cubiertos; las capturas de evidencia correspondientes se
encuentran en \texttt{docs/capturas/}.

\begin{table}[h]
    \centering
    \begin{tabular}{@{}p{5.5cm}p{2cm}p{6.5cm}@{}}
        \toprule
        \textbf{Caso} & \textbf{HTTP} & \textbf{Verificacion} \\
        \midrule
        Listado paginado, con token USER          & 200 & \texttt{meta} con \texttt{totalElements}/\texttt{totalPages} correctos \\
        Creacion con datos validos, token ADMIN    & 201 & Header \texttt{Location} presente \\
        Creacion con datos invalidos               & 400 & \texttt{errors} con \texttt{field}/\texttt{message} por cada violacion \\
        Eliminacion de producto existente, ADMIN    & 200 & \texttt{activo=false} en la fila (no se borra) \\
        Eliminacion de id inexistente                & 404 & Mensaje descriptivo \\
        Peticion sin token                          & 401 & Mensaje de no autorizado \\
        Peticion con token invalido                 & 401 & Mismo manejo que sin token \\
        POST/DELETE con token de rol USER            & 403 & Mensaje de acceso denegado \\
        Misma consulta repetida (paginacion)         & --  & Segunda vez sin log de \emph{cache miss} \\
        Alta de producto tras una lectura cacheada   & --  & Siguiente lectura vuelve a dar \emph{cache miss} \\
        \bottomrule
    \end{tabular}
    \caption{Casos de prueba ejecutados con \texttt{docs/requests.http}.}
    \label{tab:pruebas}
\end{table}

% =============================================================================
\section{Conclusiones}
% =============================================================================

El servicio implementado cubre los cinco requisitos funcionales exigidos
manteniendo una arquitectura por capas limpia, un contrato de respuesta
uniforme y verificable, y dos mecanismos transversales -- cache-aside con
Redis y seguridad JWT basada en roles -- correctamente aislados en sus
propias capas de configuracion. El uso de Docker Compose con un unico
comando de arranque, junto con un esquema de base de datos versionado
(\texttt{schema.sql}/\texttt{seed.sql}) y una coleccion de peticiones
reproducible (\texttt{requests.http}), garantiza que el proyecto pueda
verificarse de forma independiente desde una clonacion limpia del
repositorio, tal como exige el criterio de piso P3.

% =============================================================================
\bibliographystyle{plain}
\bibliography{referencias}

\end{document}
